\chapter{... que padeceu sob Pôncio Pilatos, foi crucificado, morto e sepultado; desceu aos infernos}
\chaptermark{... que padeceu sob Pôncio Pilatos}

Este capítulo aprofunda a continuação do segundo artigo do Símbolo Apostólico, professando a fé nos sofrimentos, morte, sepultamento e descida de Jesus aos infernos. Desdobramos os parágrafos correspondentes do Catecismo da Igreja Católica (CIC), explorando as controvérsias finais, a Paixão, a Cruz e o Hades. Para adultos em processo de iniciação cristã pelo RICA, esses mistérios revelam o amor redentor de Cristo, convidando à cruz pessoal no catecumenato quaresmal rumo à Páscoa.

\section{As Controvérsias Finais e a Entrada em Jerusalém (\S 571--594)}

Jesus intensifica sua missão escolhendo os Doze Apóstolos, sinal visível do novo Israel, para partilhar sua autoridade e missão. Envia-os em missão, concedendo-lhes poder sobre demônios e enfermidades, prefigurando a Igreja apostólica.

Sua pregação sobre o Reino provoca oposições crescentes: fariseus e saduceus o acusam de blasfêmia por perdoar pecados e Sabbaths violados, enquanto o povo se divide entre admiração e escândalo.

Jesus purifica o Templo, denunciando comércio profano, ato profético que precipita sua condenação. Revela-se como Filho de Davi, Rei messiânico humilde montado em jumentinho.

A entrada triunfal em Jerusalém cumpre Zc 9,13: multidões o aclamam com palmas, mas ele chora pela cidade incrédula, prevendo sua destruição.

Para adultos a partir dos 20 anos no RICA, essas controvérsias iluminam tensões modernas entre fé e mundo, preparando o rito de eleição onde os catecúmenos rejeitam o mal, abraçando a missão apostólica.

\section{A Paixão de Cristo (\S 595--623)}

A Paixão inicia no Getsêmani: Jesus, angustiado, aceita a taça do Pai em obediência filial, suando sangue. Traído por Judas, abandonado pelos discípulos, é preso como criminoso.

Julgado por Anás, Caifás e Sinédrio, falsos testemunhos o condenam por blasfêmia ao afirmar sua divindade. Pedro nega-o três vezes, cumprindo profecia, mas Jesus olha com misericórdia.

Ante Pilatos, romano cético, Jesus confessa ser Rei de verdade, não deste mundo. Flagelado, coroado de espinhos, é declarado ``Ecce Homo'', zombado como falso messias.

Carregando a Cruz ao Calvário, cai sob peso, ajudado por Simão de Cirene; Maria e Verônica o acompanham em solidariedade materna.

Na catequese para adultos não iniciados, meditar a Paixão no tempo de purificação quaresmal, com Via-Sacra, desperta compaixão e desejo de reparação, unindo sofrimentos pessoais à Cruz redentora.

\section{A Crucifixão, Morte e Sepultamento (\S 624--630)}

Na Cruz, Jesus é elevado como Rei vitorioso: ``Pai, perdoa-lhes, não sabem o que fazem.'' Crucificado entre ladrões, confia Maria a João, instituindo maternidade espiritual da Igreja.

Às trevas, clama ``Meu Deus, por que me abandonaste?'', assumindo abandono por nós; entrega o Espírito com ``Está consumado'', vitorioso sobre pecado e morte.

O centurião confessa: ``Verdadeiramente este homem era Filho de Deus''. José de Arimateia e Nicodemos sepultam-no em túmulo novo, unção apressada pelo Sabbat.

A Cruz é fonte de salvação: sacrifício perfeito, expiação universal, revelando amor trinitário até o fim.

O sepultamento afirma morte real, corpo íntegro, preparando Ressurreição.

Para catecúmenos adultos, contemplar a Cruz no scrutinium quaresmal liberta de pecados, com jejuns e esmolas imitando o sacrifício de Jesus, rumo ao Batismo pascal.

\section{Desceu aos Infernos (\S 631--637)}

``Desceu aos infernos'' significa descida às regiões dos mortos (Sheol, Hades), não aos demônios. Jesus, alma viva, abre portas do limbo dos justos, libertando Adão, profetas e patriarcas.

Esse mistério completa Paixão: pela morte, destrói morte, pregando aos espíritos em prisão (1Pd 3,19), anunciando vitória.

Os ``portões da morte não prevalecerão'': Cristo Rei penetra abismo, levando cativos ao Paraíso.

A Igreja professa isso no Credo, combatendo visões gnósticas; ícone da Ressurreição mostra-o puxando Adão e Eva.

Em catequese RICA adulta, essa verdade consola ante luto e morte, fortalecendo esperança escatológica nos exorcismos, preparando para Crisma que unge para a vitória final sobre infernos.

\section{Conclusão}

Professar a Paixão, morte e descida de Cristo revela o preço do amor salvífico. Adultos no RICA, unam-se à sua Cruz para ressurgir na Vigília Pascal.